% A.97: ordinary completed pairs and their own-coefficient projections.
\begin{statement}{A.97}{Proved}{Rigidity from the own coefficient and from one global coefficient}
Let $A,B\in\widehat{\Q^\times}$ be an ordinary completed
unit-equation pair: for every finite prime $p$ there is
$x_p\in\Q_p\setminus\{0,1\}$ with
\[
       A_p=\kappa_p(x_p),\qquad
       B_p=\kappa_p(1-x_p)
             \quad\hbox{in }\widehat{\Q_p^\times}.
\tag{A.97.1}
\]
For a coefficient prime $\ell$, put
\[
 \begin{split}
 M_\ell&=\varprojlim_n\Q^\times/(\Q^\times)^{\ell^n},\\
 M_{p,\ell}&=\varprojlim_n
                 \Q_p^\times/(\Q_p^\times)^{\ell^n}.
 \end{split}
\tag{A.97.2}
\]
Write $A_\ell\in M_\ell$ for the global coefficient
projection and $A_{p,\ell}\in M_{p,\ell}$ for its
localization; use the same notation for $B$ and rational
numbers.  These are the actual continuous homomorphisms
induced by $\Q^\times\hookrightarrow\Q_p^\times$ and by
restricting the inverse limit to the powers of $\ell$.
Fix $\lambda\in\Q\setminus\{0,1\}$.

If the own-coefficient equalities
\[
             A_{p,p}=\lambda_{p,p},\qquad
             B_{p,p}=(1-\lambda)_{p,p}
\tag{A.97.3}
\]
hold on a set of rational primes of Dirichlet density one,
then
\[
               A=\lambda,\qquad B=1-\lambda,
                    \qquad x_p=\lambda
                         \quad\hbox{for every finite }p.
\tag{A.97.4}
\]
In fact, at any single prime $p$, (A.97.3) already implies
$x_p=\lambda$.  If (A.97.4) fails, there is a nonempty
Chebotarev conjugacy class of primes, apart from finitely
many exceptions, at each of which at least one of the two
own-coefficient equalities fails.

There is also a fixed-global-coefficient criterion.
For any one coefficient prime $\ell$, the equalities
\[
                   A_\ell=\lambda_\ell,
                \qquad B_\ell=(1-\lambda)_\ell
                           \quad\hbox{in }M_\ell
\tag{A.97.5}
\]
imply (A.97.4).  Thus an ordinary completed pair cannot be
changed from a rational pair only at a proper subset of
the global coefficient primes.  This includes every
finite-primary patching of a rational pair.
The density-one assertion requires ordinary local points
only at the density-one comparison primes; under that
weaker hypothesis its conclusion is the equality of
the two full global classes $A=\lambda$, $B=1-\lambda$.
Moreover, a noncuspidal Selmer section in A.14's labeled
rational realization range, different from its realized
point section, differs from that rational point's global
Kummer coefficient pair at every prime in some nonempty
Chebotarev conjugacy class, apart from finitely many
exceptions.  In particular no finite-primary or
density-zero-primary modification gives such a section.
The fixed-global-coefficient assertion (A.97.5) retains
the all-place ordinary hypothesis (A.97.1).
No local section theorem is used in these statements;
ordinariness is the explicit hypothesis (A.97.1).
\end{statement}
\begin{proof}
We first compute the exact local kernel.  For an odd prime
$p$ the topological multiplicative decomposition is
\[
          \Q_p^\times
             =p^{\Z}\times\mu_{p-1}\times(1+p\Z_p).
\tag{A.97.6}
\]
The logarithm is an isomorphism from $1+p\Z_p$ to
$p\Z_p$.  Its intersection of all $p^n$-power subgroups
is therefore trivial.  The intersection of $p^n\Z$ in
the valuation factor is also zero.  Raising to a $p^n$th
power is an automorphism of $\mu_{p-1}$, since its order
is prime to $p$.  Consequently the exact kernel is
\[
       \ker\bigl(\Q_p^\times\longrightarrow M_{p,p}\bigr)
                         =\mu_{p-1}.
\tag{A.97.7}
\]
At two, use
$\Q_2^\times=2^{\Z}\times\{\pm1\}\times(1+4\Z_2)$.
The two-primary completion retains $\{\pm1\}$, and
the logarithm handles the last factor.  Its kernel is
trivial.  In particular (A.97.3) at two immediately
implies $x_2=\lambda$.

At an odd $p$, (A.97.1), (A.97.3), and (A.97.7) give
actual roots of unity $\zeta,\eta\in\mu_{p-1}$ with
\[
       x_p=\lambda\zeta,\qquad
       1-x_p=(1-\lambda)\eta,
       \qquad \lambda\zeta+(1-\lambda)\eta=1.
\tag{A.97.8}
\]
This is an equality in the cyclotomic number field
$L=\Q(\zeta,\eta)\subset\Q_p$.
Choose a primitive root of unity generating a cyclotomic
field containing $L$ and send it to a primitive complex
root of the same order.  Restriction gives a field
embedding $L\hookrightarrow\C$ fixing $\Q$.
Under this actual embedding $\lambda$ remains real and
$|\zeta|=|\eta|=1$.
The equality case of the triangle inequality can be
expressed without separating the three real intervals:
\[
 \begin{split}
   (1-\lambda)^2
      &=|1-\lambda\zeta|^2\\
      &=1+\lambda^2-2\lambda\operatorname{Re}(\zeta),\\
   0&=2\lambda\bigl(1-\operatorname{Re}(\zeta)\bigr).
 \end{split}
\tag{A.97.9}
\]
Since $\lambda\ne0$, we obtain
$\operatorname{Re}(\zeta)=1$, and a point of the complex
unit circle with real part one is $1$.
Then $\lambda\ne1$ and (A.97.8) give $\eta=1$ as well.
Injectivity of the field embedding proves the same
equalities in $\Q_p$.  Thus $x_p=\lambda$.

On the density-one set in the statement, both full local
classes $A_p/\lambda$ and $B_p/(1-\lambda)$ are now
trivial.  A.76's continuous localization injection gives
$A=\lambda$ and $B=1-\lambda$ globally.
At each remaining finite prime, the injectivity of
$\Q_p^\times\to\widehat{\Q_p^\times}$ from A.2 gives
$x_p=\lambda$.
For the more precise assertion, apply A.76 to either
nontrivial class $A/\lambda$ or $B/(1-\lambda)$.
It supplies a nonempty Chebotarev conjugacy class on
which that full local class is nontrivial.
The single-prime conclusion just proved shows that the
two equalities in (A.97.3) cannot both hold at any of
those primes, outside the stated finite exceptional set.

We next prove the fixed-global-coefficient assertion;
this proof does not need a density-one argument.
For each physical prime $p$, ordinary valuation gives
continuous homomorphisms
\[
 \begin{split}
   v_p &: \widehat{\Q^\times}\longrightarrow\widehat\Z,\\
   v_{p,\ell} &: M_\ell\longrightarrow\Z_\ell.
 \end{split}
\tag{A.97.10}
\]
They are induced at each finite level by the ordinary
integer valuation modulo the relevant exponent.  The
square consisting of these two maps and the coefficient
projections commutes, and localization commutes with it
as well.  By (A.97.1), $v_p(A)$ is the diagonal image
in $\widehat\Z$ of the ordinary integer $v_p(x_p)$.
Thus (A.97.5) gives
\[
        v_p(x_p)=v_p(\lambda)\quad\hbox{in }\Z_\ell,
       \qquad
        v_p(1-x_p)=v_p(1-\lambda)\quad\hbox{in }\Z_\ell.
\tag{A.97.11}
\]
The map $\Z\hookrightarrow\Z_\ell$ is injective, so
these are equalities of ordinary integers for every $p$.
In particular all full valuation coordinates of
$A/\lambda$ and $B/(1-\lambda)$ vanish.

The kernel of the actual valuation map
\[
          \widehat{\Q^\times}\longrightarrow
                     \prod_p\widehat\Z,
                 \qquad C\longmapsto(v_p(C))_p
\tag{A.97.12}
\]
is precisely $\{\pm1\}$.
Here is a finite-level proof, which does not assert
surjectivity onto the displayed product.
Choose a rational representative $c_n$ of a class
modulo $n$ in the kernel.  Each prime valuation of $c_n$
is divisible by $n$, so rational prime factorization
gives $c_n=\epsilon_n d_n^n$ with
$\epsilon_n\in\{\pm1\}$ and $d_n\in\Q^\times$.
For odd $n$ the sign is itself an $n$th power.
For even $n$ the sign is determined by the common
coordinate modulo two.  Compatibility therefore makes
the entire inverse-limit class one fixed sign.
Conversely both signs have zero valuations.

It follows that
$A=\epsilon\lambda$ and $B=\delta(1-\lambda)$ for
some signs $\epsilon,\delta$.
Full local Kummer injectivity and (A.97.1) imply, at any
one finite prime,
\[
                   \epsilon\lambda+
                           \delta(1-\lambda)=1
                         \quad\hbox{in }\Q.
\tag{A.97.13}
\]
The signs $(+,-)$ and $(-,+)$ force respectively
$\lambda=1$ and $\lambda=0$; the signs $(-,-)$ give
$-1=1$.  All are excluded.  Hence both signs are positive,
which proves (A.97.4).

Finally the prime-power coefficient map
\[
             \widehat{\Q^\times}
                         \longrightarrow\prod_\ell M_\ell
\tag{A.97.14}
\]
is an isomorphism of topological abelian groups:
at a finite exponent this is the Chinese remainder
isomorphism for the abelian group $\Q^\times$, and
passing to the indicated inverse limits gives the map
and its continuous inverse.  If a modification leaves
even one global coefficient component of both entries
equal to the rational pair, (A.97.5) applies.
For example a modification supported at one primary
idempotent $e_\ell\in\widehat\Z$ has the form
\[
              A=\lambda C^{e_\ell},\qquad
              B=(1-\lambda)D^{e_\ell},
\tag{A.97.15}
\]
where exponentiation is defined in each finite quotient
and then in the inverse limit.  Every coefficient
$\ell'\ne\ell$ is unchanged, so an ordinary completed
pair of this form is the rational pair itself.
The same argument covers any proper set of modified
coefficient primes, not only a single one or a finite set.
\end{proof}

\paragraph{Ordinariness on a density-one set and Selmer sections.}
The density-one assertion has the following precise weaker
hypothesis.  Let $A,B$ be arbitrary completed global classes,
and suppose that on one density-one set of finite primes
there are ordinary points satisfying both (A.97.1) and
(A.97.3).  Then $A=\lambda$ and $B=1-\lambda$ globally.
The proof is exactly (A.97.6)--(A.97.9) followed by A.76;
it uses no local point at an omitted prime.  At any further
prime where an ordinary point does satisfy (A.97.1), that
point is consequently $\lambda$.

In particular, let $s$ be a noncuspidal Selmer section in
the labeled rational realization range of A.14, and fix
its realized parameter $\lambda$.
If the global coefficient pairs
$(A_{s,\ell},B_{s,\ell})$ equal
$(\lambda_\ell,(1-\lambda)_\ell)$ outside a set of
coefficient primes of Dirichlet density zero, then
$s=s_\lambda$.
Indeed, A.37 restricts the possible finite locally
cuspidal places to the finitely many potentially
multiplicative primes of the realized curve.
At every other physical prime the section has an ordinary
local point by the Selmer hypothesis.
Intersecting these places with the unchanged coefficient
primes gives the density-one set required above.
Hence the full global Kummer pair is rational.

For clarity, this also excludes every finite locally
cuspidal place directly.  The actual rational functions
regular with value one at the three cusps are
\[
       1-t\quad\hbox{at }0,\qquad
       t\quad\hbox{at }1,\qquad
       \frac{t}{t-1}\quad\hbox{at }\infty.
\tag{A.97.16}
\]
Their pullbacks in the respective cusp packets have
Kummer class one.  For the rational pair these classes
are respectively $1-\lambda$, $\lambda$, and
$-\lambda/(1-\lambda)$.  Local Kummer injectivity at
a finite prime would force the corresponding rational
number to equal one, and each of the three equalities
contradicts $\lambda\ne0,1$.
All finite local centres are therefore the point
$\lambda$, and finite-support recognition, equivalently
the no-local-cusp case of A.37, gives $s=s_\lambda$.
This argument uses prescribed Selmer local data, rather
than an additional use of the local section theorem.

Equivalently, any different Selmer section in this
realization range must differ from the rational Kummer
pair in its global coefficient components at every prime
of some nonempty Chebotarev conjugacy class, apart from
finitely many exceptions.  To see this, apply A.76 to a
nontrivial full Kummer ratio, omit the finite potentially
multiplicative set, and use the single-prime implication
(A.97.3)$\Rightarrow x_p=\lambda$.
In particular, finite-primary and density-zero-primary
patches do not produce such a section.
The separate fixed-global-coefficient proof
(A.97.10)--(A.97.13) still requires ordinary local
valuations at every finite place; it is not asserted for
arbitrary cuspidal local data.

\paragraph{The exact remaining comparison in A7.}
A.56 gives the continuous homomorphism from the geometric
abelianization $T=\widehat\Z(1)^2$ to the abelianization
of the full geometric congruence image:
\[
        T\longrightarrow H^{\mathrm{ab}},\qquad
         (a,b)\longmapsto
              (a\bmod8,\ b\bmod8,\ a+b\bmod3).
\tag{A.97.17}
\]
The equality of full $H$-sections induces the comparisons
under this actual map.  The coefficient equalities
(A.97.3) or (A.97.5) are the additional, explicitly
stated data used in A.97.  The remaining recognition
problem is to use the full $H$-comparison and its
ordinary or cuspidal local data to prove rationality;
no passage from (A.97.17) to a full coefficient equality
is assumed here.
